% Introduction
\section{Introduction}

The domestic pig (\textit{Sus scrofa}) is an indispensable model for biomedical research due to its anatomical and physiological homology to humans \cite{swindle2012,lunney2021}. Yet, a critical gap limits its translational utility: the lack of a standardized, molecular definition of development. Current staging relies on morphological criteria or chronological age, which are often poor proxies for biological maturity. Littermates can develop asynchronously, and distinct organs mature at different rates effectively decoupling "whole-body" age from tissue-specific functional competence. This imprecision confounds experimental reproducibility and obscures subtle developmental phenotypes, particularly in studies of pediatric disease or regenerative medicine \cite{bendus2006,orrahilly2010}.

Molecular clocks offer a solution. Transcriptomic profiles provide a quantitative readout of cellular state, enabling precise "biological aging" that transcends morphological variation \cite{schachtschneider2021,zhengh2023}. While such molecular clocks are established in humans and mice, a comprehensive, tissue-resolved developmental atlas has been missing for the pig.

Here, we present the first calibrated transcriptomic atlas of porcine development, integrating data across five major tissues (Brain, Liver, Muscle, Lung, and Blood) with sufficient sample sizes for machine learning analysis. Unlike previous descriptive studies, we build a quantitative, predictive framework that infers developmental stage with high accuracy. Most significantly, we demonstrate that this porcine molecular clock captures evolutionarily conserved developmental programs in muscle tissue: our cross-species analysis reveals a striking quantitative correlation with human developmental trajectories (Pearson R = 0.62), identifying shared regulatory modules that orchestrate the transition from infant to adult neuromuscular maturity. We highlight conserved roles for \textit{SORCS2} and \textit{FGFRL1} in establishing synaptic connectivity and fiber specialization in both species. We note that cross-species validation is currently limited to muscle tissue due to availability of matched human developmental data; extension to other tissues represents an important future direction. This work establishes a rigorous new standard for staging in the pig model, directly bridging the gap between porcine experiments and human clinical translation, with particular strength for muscle-related applications.